\section{ПРАКТИЧЕСКАЯ АПРОБАЦИЯ, ОГРАНИЧЕНИЯ И НАПРАВЛЕНИЯ РАЗВИТИЯ}

После проектирования, реализации и экспериментальной оценки необходимо определить, какую практическую задачу решает разработанная платформа, где проходят границы текущего MVP и какие направления развития являются наиболее важными. Данная глава связывает технические результаты предыдущих разделов с эксплуатационным смыслом работы: снижением ручной нагрузки на инженеров и повышением наблюдаемости состояния инфраструктуры.

\subsection{Практическая апробация разработанной платформы}

Практическая апробация выполнялась на сценарии управления контейнерными workload, который отражает основную идею работы. В архитектурном описании задается управляемый хост и два требуемых компонента: \texttt{node\_exporter} и \texttt{cadvisor}. ParserSvc преобразует это описание в desired state, где для каждого workload указаны имя, признак включения, режим развертывания, image и port. InventorySvc формирует actual state, получая сведения о фактически наблюдаемых контейнерных workload. DriftDetectorSvc сравнивает состояния и при подтвержденном расхождении отправляет reconcile-команду в ReconcilerSvc.

Сценарий выбран не случайно. \texttt{node\_exporter} и \texttt{cadvisor} относятся к инфраструктурным компонентам наблюдаемости. Такие workload часто должны присутствовать на управляемых хостах, но их отсутствие не всегда сразу заметно прикладным командам. Если агент или контейнер наблюдаемости пропал, эксплуатационная команда может потерять часть метрик и обнаружить проблему уже после того, как понадобится диагностика инцидента. Поэтому автоматическое обнаружение отсутствующего компонента и запуск восстановления имеют понятную практическую ценность.

В ходе апробации подтверждено, что платформа реализует полный путь данных. Desired state формируется из Structurizr JSON, а не задается вручную внутри DriftDetectorSvc. Actual state поступает через InventorySvc и содержит не только список workload, но и metadata о полноте данных. DriftDetectorSvc учитывает эту полноту и не отправляет reconcile-команды при недостоверном actual state, кроме специально предусмотренного bootstrap-сценария для \texttt{cadvisor}. ReconcilerSvc принимает команду асинхронно и передает выполнение operator-слою.

Отдельный результат апробации связан с partial-семантикой. В инфраструктурной автоматизации ошибочное действие может быть опаснее бездействия. Поэтому платформа должна уметь признавать неполноту данных и не превращать сетевую ошибку или недоступность источника в автоматический reconcile. Проверенный сценарий partial data подтвердил, что данное требование учтено: недостоверные данные приводят к предупреждениям и partial-статусу, а не к неконтролируемому восстановлению.

Таким образом, практическая апробация подтвердила реализуемость основной идеи: Architecture-as-Code может использоваться как источник desired state для автоматизированного управления инфраструктурными workload, если дополнить его inventory-слоем, detector-ами и ограниченным reconcile-контуром.

\subsection{Практическая значимость платформы для SRE/DevOps-команд}

Для SRE- и DevOps-команд ключевая ценность платформы состоит в снижении ручной нагрузки при контроле состояния инфраструктуры. Без такого инструмента инженер вынужден отдельно поддерживать архитектурное описание, inventory, playbook и процедуры проверки. При изменении инфраструктуры требуется вручную понимать, какие хосты должны содержать какие компоненты, какие из них фактически присутствуют и какие действия нужно выполнить для восстановления.

Платформа предлагает другой подход. Архитектурное описание становится источником desired state, inventory-сервис предоставляет actual state, а drift detection выполняет сопоставление автоматически. Это снижает число ручных операций и уменьшает вероятность того, что расхождение будет пропущено из-за человеческого фактора.

Практическая значимость проявляется в нескольких аспектах:
\begin{enumerate}
    \item повышение наблюдаемости расхождений между архитектурным описанием и фактической инфраструктурой;
    \item уменьшение времени обнаружения отсутствующих инфраструктурных workload;
    \item формализация сценариев восстановления через operator-слой;
    \item снижение зависимости от ручного запуска отдельных playbook;
    \item возможность дальнейшего расширения платформы без переработки базовой архитектуры.
\end{enumerate}

Для команд, которые уже используют архитектурные описания в виде кода, платформа создает дополнительную ценность: архитектурная модель становится не только документацией, но и входом для автоматизированного контроля. Это повышает актуальность архитектурных артефактов. Если описание используется в рабочем процессе платформы, у команды появляется практическая мотивация поддерживать его в согласованном состоянии.

С точки зрения эксплуатации платформа может стать основой внутреннего control plane для host-level workload. В текущем виде она решает ограниченный сценарий, но архитектура допускает добавление новых источников actual state, новых detector-ов и новых operator-ов. Благодаря этому решение может развиваться постепенно, начиная с наиболее типовых инфраструктурных компонентов и переходя к более сложным сценариям.

\subsection{Ограничения текущего MVP}

Несмотря на подтвержденную работоспособность, текущая версия платформы имеет ряд ограничений. Эти ограничения важно фиксировать явно, поскольку они определяют область корректного применения результата и направления дальнейшего развития.

Первое ограничение связано с источником desired state. ParserSvc работает со Structurizr JSON и не выполняет разбор Structurizr DSL напрямую. Это означает, что подготовка JSON-представления должна выполняться до передачи данных в сервис. Такое решение упростило MVP, но в дальнейшем может быть дополнено автоматизированным pipeline подготовки манифестов или встроенной поддержкой дополнительных форматов.

Второе ограничение связано с моделью actual state. В текущей реализации реально интегрирован cAdvisor как источник сведений о контейнерных workload. NetBox и другие системы инвентаризации обозначены как направления расширения, но не участвуют в формировании actual state в MVP. Поэтому платформа пока не решает задачу полноценной инвентаризации всей инфраструктуры.

Третье ограничение состоит в отсутствии персистентного хранения snapshot и истории. Desired state, actual state и cooldown-состояние хранятся в памяти сервисов. Это удобно для проверки архитектурной идеи, но недостаточно для зрелой эксплуатации. После перезапуска сервисов история detection cycle, предыдущие snapshot и сведения о подавленных reconcile-командах теряются.

Четвертое ограничение относится к drift-модели. В MVP drift определяется как отсутствие поддержанного контейнерного workload при наличии соответствующего требования в desired state. Более сложные типы drift пока не реализованы.

Пятое ограничение связано с ReconcilerSvc. Сервис возвращает \texttt{202 Accepted} после постановки задачи в очередь, но не предоставляет API получения статуса выполнения. Результат Ansible Runner фиксируется в логах, однако отсутствует persisted job state, история запусков, retry policy и dead-letter-механизм.

Шестое ограничение касается безопасности. В текущей версии не реализованы полноценные механизмы аутентификации и авторизации между сервисами, audit trail и политики доступа к reconcile-действиям. Поскольку платформа потенциально выполняет изменения на управляемых хостах, развитие безопасности является обязательным условием дальнейшего использования.

Седьмое ограничение связано с масштабированием и отказоустойчивостью. Экспериментальная оценка подтверждает работоспособность текущей архитектуры, но не является промышленным SLA. Не проверены сценарии длительной работы, отказов отдельных сервисов, восстановления после перезапуска, горизонтального масштабирования и конкурентного выполнения нескольких экземпляров одного сервиса.

Таким образом, текущий MVP следует рассматривать как доказательство реализуемости архитектурного подхода, а не как завершенную промышленную платформу. Его ценность состоит в том, что он подтверждает правильность декомпозиции и работоспособность базового control loop.

\subsection{Направления дальнейшего развития}

Дальнейшее развитие платформы можно разделить на несколько направлений.

Первое направление --- расширение источников actual state. Наиболее естественным следующим шагом является интеграция NetBox как источника сведений о хостах, ролях, окружениях, сетевых атрибутах и принадлежности ресурсов. Это позволит уменьшить зависимость от bootstrap-файлов и приблизить InventorySvc к полноценной модели инфраструктурного inventory. Также могут быть добавлены другие источники: CMDB, агентские системы, Prometheus-метрики или специализированные API платформенных команд.

Второе направление --- persistence. Необходимо добавить персистентное хранение snapshot desired state и actual state, истории detection cycle, найденных drift и отправленных reconcile-команд. Это позволит анализировать динамику состояния инфраструктуры, восстанавливать контекст после перезапуска сервисов и строить отчеты по качеству управления инфраструктурой.

Третье направление --- развитие метрик и наблюдаемости. Для DriftDetectorSvc полезны метрики по длительности стадий цикла, количеству desired/actual-хостов, числу partial-циклов, найденным drift, отправленным и подавленным reconcile-командам. Для ReconcilerSvc важны queue depth, число активных worker-ов, длительность выполнения задач, доля успешных и ошибочных запусков. Такие метрики позволят перейти от разовых измерений к постоянному мониторингу качества работы платформы.

Четвертое направление --- расширенная drift-модель. Помимо отсутствия workload, необходимо проверять версию образа, параметры запуска, порт, режим развертывания, состояние здоровья и, при необходимости, конфигурационные файлы. Это потребует более богатой модели desired state, расширения actual state и реализации новых detector-ов.

Пятое направление --- новые operators. Текущий набор ограничен \texttt{node\_exporter} и \texttt{cadvisor}. В дальнейшем можно добавить operators для других компонентов наблюдаемости и инфраструктурных сервисов: blackbox exporter, alertmanager, log collector, node-level security agents, backup agents и других типовых workload. При этом важно сохранять общий контракт: detector обнаруживает drift, operator строит execution plan, executor выполняет восстановление.

Шестое направление --- развитие ReconcilerSvc. Для зрелой эксплуатации нужны persisted job state, API статуса задачи, retry policy с backoff, dead-letter queue, приоритеты reconcile-команд и ограничения по одновременному воздействию на один хост. Это позволит сделать восстановительные действия более управляемыми и безопасными.

Седьмое направление --- безопасность. Платформа должна поддерживать аутентификацию сервисов, авторизацию действий, audit trail, контроль прав на запуск reconcile, защиту секретов и политику допуска operator-ов. Без такого слоя автоматизированная система управления инфраструктурой не может считаться готовой к использованию в критичных средах.

Восьмое направление --- развитие Architecture-as-Code-контура. В перспективе можно добавить проверку схемы архитектурных манифестов, расширенные правила валидации, поддержку нескольких архитектурных источников, версионирование desired state и сравнение изменений между версиями архитектурного описания. Это усилит связь между архитектурной документацией и эксплуатационным состоянием.

Перечисленные направления не требуют отказа от текущей архитектуры. Напротив, они естественно продолжают заложенную декомпозицию. ParserSvc может развиваться в сторону более богатого desired state, InventorySvc --- в сторону расширенного actual state, DriftDetectorSvc --- в сторону новых detector-ов и политик, ReconcilerSvc --- в сторону надежного job management и безопасного исполнения.

\subsection{Выводы по практической апробации}

Практическая апробация показала, что разработанная платформа решает ключевую задачу MVP: связывает Architecture-as-Code-описание с автоматизированным контролем фактического состояния инфраструктуры. Сценарий с контейнерными workload подтвердил, что desired state может быть построен из Structurizr JSON, actual state может быть получен через inventory-слой, drift может быть обнаружен отдельным сервисом, а восстановление может быть выполнено через ReconcilerSvc.

Для SRE- и DevOps-команд такой подход полезен тем, что снижает зависимость от ручных проверок и отдельных запусков playbook. Платформа формирует основу для более системного управления инфраструктурой: архитектурное описание, inventory, detection и reconcile становятся частями одного процесса.

В то же время текущая реализация остается MVP. Она ограничена двумя workload, in-memory состоянием, cAdvisor как основным источником actual state и отсутствием полноценной истории выполнения reconcile-задач. Эти ограничения не снижают ценность полученного результата, но задают границы его применения.

Дальнейшее развитие должно быть направлено на расширение источников inventory, добавление persistence, развитие метрик, усложнение drift-модели, увеличение набора operator-ов и усиление безопасности. При реализации этих направлений платформа может эволюционировать из исследовательского MVP в полноценный внутренний control plane для управления инфраструктурными workload.

Таким образом, цель практической апробации достигнута: подтверждена работоспособность основной идеи автоматизированной платформы управления инфраструктурой с поддержкой подхода Architecture-as-Code, выявлены ограничения текущей реализации и сформирован реалистичный план дальнейшего развития.
